\documentclass[10pt,a4paper]{amsart}
\usepackage[margin=19mm]{geometry}
\usepackage{amsmath,amssymb,mathtools}
\usepackage[scr=rsfs,cal=euler]{mathalfa}
\usepackage{xfrac}
\usepackage[unicode,hidelinks,bookmarks=false]{hyperref}
\newcommand{\ie}{i.e.\ }
\newcommand{\cf}{cf.\ }
\newcommand{\resp}{respectively\ }
\newcommand{\Subsection}{\subsection}
\providecommand{\nobreakdash}{\nobreak\hbox{-}}
\setlength{\emergencystretch}{2em}
\multlinegap0pt
\usepackage{amssymb}
\usepackage{bm}
\newtheorem{lemma}{Lemma}
\title{On the Maximal Length of MDS Elliptic Codes}
\author{Haojie Chen, Chuangqiang Hu\(\dagger\) , Junjie Huang, Chang-An Zhao}
\date{Source: arXiv:2605.29439v2 (CC BY 4.0)}
\begin{document}
\maketitle

\noindent\emph{Source and licence.} On the Maximal Length of MDS Elliptic Codes. Version arXiv:2605.29439v2 (CC BY 4.0), distributed by arXiv under the Creative Commons Attribution 4.0 International licence, http://creativecommons.org/licenses/by/4.0/, as recorded in the arXiv licence metadata for this exact version and shown on the abstract page https://arxiv.org/abs/2605.29439v2. Full licence notice: \texttt{licenses/arxiv-2605.29439-CC-BY-4.0.txt}.\par\smallskip

\noindent\textit{Editorial scope.} Counted unit: the article's lemma lemma:order\_two (source0.tex lines 426--428). The article's complete original proof of it is delivered with it (source0.tex lines 430--446, one \texttt{proof} environment of 17 lines). No mathematics is written, repaired or re-derived: the statement and the proof are the article's own lines, in the article's own notation, and no retained line is substituted or re-flowed.  No further source-local context is needed: every cross-reference of every retained line resolves inside the two delivered spans.  Nothing was omitted from any retained span, so the package carries no omission record.  No retained line contains a citation, so the wrapper carries no bibliography and no bibliography entry.  No retained line contains a figure environment, an image-inclusion command or drawing code, so no image is delivered and the package declares no asset dependency.  Editorial formatting only, in addition: an AMS \texttt{amsart} preamble that restates the article's own declarations and macros used by the retained lines, namely the environments \texttt{lemma} on separate counters, together with the standard packages those lines need.  The retained environments print in this document as Lemma 1 in order of appearance, whereas the article prints the same items with the numbering of its own full document.  The complete native source of the article as distributed in the arXiv e-print for this exact version is bundled unchanged as \texttt{sources/201609/source0.tex}.
\begin{lemma}\label{lemma:order_two}
    Let $\mathbf{G}$ be a finite abelian group and let $\mathbf{G}[2] = \{g \in \mathbf{G} \mid 2g = 0\}$ be its $2$-torsion subgroup. Then the number of subgroups $H \le \mathbf{G}$ of index $2$ equals $|\mathbf{G}[2]| - 1$.
\end{lemma}

\begin{proof}
    Consider the homomorphism $\varphi: \mathbf{G} \to \mathbf{G}$ given by $\varphi(g) = 2g$. Its image is $2\mathbf{G} = \{2g \mid g \in \mathbf{G}\}$ and its kernel is $\mathbf{G}[2]$. Thus we have an exact sequence
    \[
    0 \longrightarrow \mathbf{G}[2] \longrightarrow \mathbf{G} \stackrel{\varphi}{\longrightarrow} \mathbf{G} \longrightarrow \mathbf{G}/2\mathbf{G} \longrightarrow 0,
    \]
    where the last map is the canonical projection. Since $\mathbf{G}$ is finite, we get
    \[
    |\mathbf{G}/2\mathbf{G}| = |\mathbf{G}[2]|.
    \]
    Moreover, $\mathbf{G}/2\mathbf{G}$ is an elementary abelian $2$-group (since $2(g+2\mathbf{G})=2g+2\mathbf{G}=0$ in the quotient), and the same holds for $\mathbf{G}[2]$. Therefore both are $\mathbb{F}_2$-vector spaces, and they have the same dimension, say $r$.

    Every subgroup $H \le \mathbf{G}$ of index $2$ is the kernel of a non‑trivial homomorphism $\chi: \mathbf{G} \to \{\pm 1\} \cong \mathbb{F}_2$ (viewed additively). Conversely, the kernel of any such non‑trivial character is a subgroup of index $2$. Distinct characters give distinct kernels because a non‑trivial character is uniquely determined by its kernel (the target has only two elements). Hence the number of index‑$2$ subgroups equals the number of non‑trivial elements of the dual group $\widehat{\mathbf{G}} = \operatorname{Hom}(\mathbf{G}, \mathbb{F}_2)$.

    For any finite abelian group, every homomorphism $\mathbf{G} \to \mathbb{F}_2$ factors through $\mathbf{G}/2\mathbf{G}$ because $2\mathbf{G}$ is contained in the kernel (as $\mathbb{F}_2$ is of exponent $2$). Thus $\widehat{\mathbf{G}} \cong \operatorname{Hom}(\mathbf{G}/2\mathbf{G}, \mathbb{F}_2)$. Since $\mathbf{G}/2\mathbf{G}$ is an $\mathbb{F}_2$-vector space of dimension $r$, its dual space is isomorphic to $\mathbf{G}/2\mathbf{G}$ itself, and therefore $|\widehat{\mathbf{G}}| = |\mathbf{G}/2\mathbf{G}| = 2^r = |\mathbf{G}[2]|$.

    Consequently, the number of non‑trivial characters, and hence the number of subgroups of index $2$, is $|\mathbf{G}[2]| - 1$.
\end{proof}

\end{document}
